\documentclass[11pt,a4paper]{ctexart}
\usepackage{amsmath,amssymb,bm}
\usepackage{booktabs}
\usepackage{longtable}
\usepackage[margin=2.3cm]{geometry}
\usepackage[hidelinks]{hyperref}
\usepackage{enumitem}
\usepackage{caption}
\usepackage{array}
\captionsetup{font=small,labelfont=bf}
\linespread{1.18}
\setlist{itemsep=2pt,topsep=3pt}

\newcommand{\R}{\mathbb{R}}
\newcommand{\norm}[1]{\left\lVert #1 \right\rVert}

\title{\bfseries 质量不平衡与不对中的等效扰动建模\\[2mm]
\large ——为什么故障可以写成正弦叠加：文献综述}
\author{RCESO 项目组}
\date{2026-11-24}

\begin{document}
\maketitle

\begin{abstract}
\noindent
本文回答一个具体的建模缺口：在等效输入扰动（EID）框架中，扰动集被取为有限个正弦之和
$\Phi=\{p_i(t)\sin(\omega_i t+\varphi_i)\}$，但原始文献并未论证\textbf{质量不平衡}与\textbf{不对中}为何落在该集合内。

综述表明，这一结论由两步构成：
第一，旋转机械的故障具有\textbf{旋转周期性}，因而其激励是转角的周期函数——
不平衡的离心力为 $me\omega^2e^{j\omega t}$，严格落在 $1\times$；
不对中使联轴器刚度呈现 $K_{\rm SCS}+K_{\rm ACS}(\theta)$ 的时变特性，参数激励使响应必然含 $1\times,2\times,\dots$ 谐波。
第二，任何周期函数都可按傅里叶级数展开为有限（截断后）正弦之和。
因此\textbf{正弦形式是数学结论而非建模假设}，物理假设只有“故障在旋转中周期复现”这一条。

本文以顶级期刊文献为主干（IEEE TIE、MSSP、JSV、Cambridge 专著），
并标注每篇的分量、被引次数与在论证链中的作用；
对无法获取全文的订阅文献，另列开放获取（CC-BY）的补充文献，已随文归档。
\end{abstract}

\begin{quote}
\small\itshape
说明：标注 \(\dagger\) 的文献为订阅制，本文仅据其题录、摘要与可获取片段归纳，
未通读全文；引用具体论断前请核对原文。
\end{quote}

\tableofcontents
\newpage

\section{问题的提出}

EID 方法把一切未知扰动折算为作用在控制输入端的等效扰动，
其理论依据是 She 等给出的存在性引理：
在一定假设下，控制输入通道上\textbf{总存在}一个 EID，
且它产生的输出与原扰动产生的输出完全相同。

然而该引理只保证“存在”，并未回答\textbf{具体形态}问题。
原文直接取扰动集为
\begin{equation}
\Phi=\left\{\,p_i(t)\sin(\omega_i t+\varphi_i)\,\right\},\qquad i=0,\dots,n,\ n<\infty,
\end{equation}
即有限个正弦（幅值可含多项式慢变）之和。

对旋转机械故障诊断而言，这个取法是否成立，需要物理建模来回答。
本文的结论是：\textbf{成立，且由两条独立的论证共同支撑}。

\section{论证链概览}

\begin{center}
\fbox{\parbox{0.82\textwidth}{\centering\vspace{2pt}
\textbf{① 故障几何具有旋转周期性}\\[1mm]
$\downarrow$ \quad 物理建模（\S\ref{sec:unb}、\S\ref{sec:mis}）\\[1mm]
\textbf{② 激励是转角 $\theta=\omega_{\rm rot}t$ 的周期函数}\\[1mm]
$\downarrow$ \quad 傅里叶级数（\S\ref{sec:fourier}）\\[1mm]
\textbf{③} $\displaystyle\sum_{i}p_i\sin(i\,\omega_{\rm rot}t+\varphi_i)\in\Phi$
\vspace{2pt}}}
\end{center}

关键点在于：\textbf{物理假设只有一条}——故障在旋转中周期复现；
正弦形式是傅里叶分解的必然结果。
需要物理建模的只是\textbf{含哪些阶次}。

\section{质量不平衡：严格的一倍频}
\label{sec:unb}

\subsection{物理推导}

设转子偏心质量 $m$、偏心距 $e$、角速度 $\omega$。
旋转坐标系中离心力大小为
\begin{equation}
F_u=m\,e\,\omega^2,
\end{equation}
方向固定于偏心方向。
在\textbf{固定坐标系}中，该力随转子一起旋转，故为旋转矢量
\begin{equation}
\vec F_u(t)=m\,e\,\omega^2\,e^{j\omega t}.
\end{equation}

由此得到三条结论：
\begin{enumerate}[label=(\alph*)]
\item \textbf{只有 $1\times$ 一个阶次}，不存在高次谐波；
\item 幅值 $\propto\omega^2$，随转速平方增长；
\item 相位锁定转子相位。
\end{enumerate}

这正是振动监测中“$1\times$ 幅值上升即为不平衡”这一判据的物理来源。

\subsection{文献}

\begin{table}[h]
\centering
\caption{不平衡建模的依据}
\begin{tabular}{@{}p{7.2cm}ll@{}}
\toprule
文献 & 类型 & 被引 \\
\midrule
Friswell, Penny, Garvey, Lees, \emph{Dynamics of Rotating Machines},
Cambridge University Press, 2010. DOI 10.1017/CBO9780511780509
& \textbf{权威专著}（旋转机械标准教材） & \textbf{422} \\
\bottomrule
\end{tabular}
\end{table}

\subsection{重力项：水平轴上 $1\times$ 转矩的来源与文献依据}
\label{sec:gravity}

前两小节的离心力是\textbf{径向力}，对转轴的力矩为零。
但在\textbf{水平轴}机器上，偏移质心还会受到重力作用，而重力方向与质心位置矢量\textbf{不平行}，
其力矩非零，且正是 $1\times$ 转矩的主要来源。

\subsubsection{推导}

设转轴水平，重力沿竖直向下。取转角 $\theta$ 自竖直向下量起，质心位置
\begin{equation}
\vec r_c=e\left(-\sin\theta,\ -\cos\theta\right),\qquad
\vec F_g=\left(0,\ -Mg\right).
\end{equation}
绕转轴的力矩为
\begin{equation}
T_g=\left(\vec r_c\times\vec F_g\right)_z
=r_xF_{gy}-r_yF_{gx}
=(-e\sin\theta)(-Mg)=M\,g\,e\sin\theta .
\end{equation}

该式与\textbf{复摆}（physical pendulum）的重力矩完全同形：重力恒定，而力臂随转角按 $\sin\theta$ 变化，
故为严格的 $1\times$ 分量，与转速无关。

\subsubsection{两个必须说明的性质}

\textbf{（1）保守性。} 重力矩在一整转上的积分为零：
\begin{equation}
\int_0^{2\pi}Mge\sin\theta\,\mathrm{d}\theta=0 .
\end{equation}
因此它只产生\textbf{转速脉动}，不贡献\textbf{平均负载}。
若控制目标是抑制平均负载，此项不应计入；若目标是抑制纹波，则它是主项。

\textbf{（2）与摩擦调制项的交叉。} 轴承摩擦调制项为 $\mu_b r_bMe\omega^2$，与 $\omega^2$ 成正比，
而重力项与 $\omega$ 无关。两者之比
\begin{equation}
\frac{T_{\rm fric}}{T_{\rm grav}}=\frac{\mu_b r_b\,\omega^2}{g},
\end{equation}
交叉角频率
\begin{equation}
\omega^{*}=\sqrt{\frac{g}{\mu_b r_b}} .
\end{equation}
取 $\mu_b=0.01$、$r_b=0.01$ m，得 $\omega^{*}\approx313$ rad/s（约 3000 r/min）。
故\textbf{低速段重力项主导，高速段摩擦调制项主导}。

\subsubsection{文献依据}

重力力矩是刚体静力学的基本结论，无需专门文献“提出”；
需要文献支撑的是\textbf{把重力项显式写入转动系统运动方程}这一做法，
以及\textbf{机械故障经此通道转化为 $1\times$ 转矩}的建模实践。
下列文献均属此类，且以顶级期刊为主。

\begin{table}[h]
\centering
\caption{显式包含重力项的转动系统建模文献}
\small
\begin{tabular}{@{}p{7.3cm}p{3.2cm}r@{}}
\toprule
文献 & 出处与分量 & 被引 \\
\midrule
Kallesøe (2007), \emph{Equations of motion for a rotor blade, \textbf{including gravity}, pitch action and rotor speed variations}
& \textbf{Wind Energy} 10:209--230（风能顶刊，Q1） & 56 \\
Yang, Huang et al. (2025), \emph{Investigation on vibration characteristics of a rotor system with \textbf{combined misalignment and eccentric faults}}
& \textbf{MSSP} 240:113332（\textbf{机械信号处理顶刊}） & 6 \\
Chung \& Jang (2003), \emph{Dynamic response and stability analysis of an \textbf{automatic ball balancer} for a flexible rotor}
& \textbf{JSV} 259:31--43 & 49 \\
Haidar \& Palacios (2021), \emph{Modified ball-type automatic balancer for rotating shafts}
& \textbf{JSV} 496:115927 & 15 \\
Rezaee \& Fathi (2015), \emph{Improving the working performance of automatic ball balancer}
& \textbf{JSV} 358:375--391 & 10 \\
Bhattacharyya \& Dutt (1997), \emph{Unbalance response and stability analysis of \textbf{horizontal rotor} systems mounted on nonlinear rolling element bearings}
& \textbf{ASME J. Vibration and Acoustics} 119:539--544 & 10 \\
Song \& Liu (2024), \emph{Research on unbalanced characteristics of axial-varying eccentric rotor}
& \textbf{MSSP} 220:111647 & 13 \\
\bottomrule
\end{tabular}
\end{table}

\noindent
\textbf{说明}：以重力力矩作为\textbf{电机驱动负载扰动}的文献在公开检索中较少，
原因是该量级仅 $10^{-4}$ N$\cdot$m，工程上常被忽略。
因此本文将其定位为\textbf{建模完整性项}：其物理来源清楚（式 (7)），
其方程形式有上列顶刊文献的建模先例，其保守性已明确声明。
这一处理方式比强行引用一篇“提出该公式”的文献更为诚实。

\section{不对中：时变刚度与谐波生成}
\label{sec:mis}

\subsection{核心机制：静态刚度 + 时变附加刚度}

不对中之所以产生 $1\times$、$2\times$ 乃至更高阶谐波，目前最清晰的物理解释来自 Lees (2007)：
存在不对中时，联轴器刚度可分解为
\begin{equation}
K(\theta)=K_{\rm SCS}+K_{\rm ACS}(\theta),
\end{equation}
其中 $K_{\rm SCS}$ 为\textbf{静态联轴器刚度}（不随转角变化），
$K_{\rm ACS}(\theta)$ 为\textbf{时变附加刚度}，随转角周期变化。

$K_{\rm ACS}$ 的周期由不对中类型决定：

\begin{table}[h]
\centering
\caption{不对中类型与主导阶次}
\begin{tabular}{@{}lcl@{}}
\toprule
类型 & 每转刚度变化周期 & 主导阶次与物理图像 \\
\midrule
平行不对中 & 2 次（拉--压交替） & \textbf{$2\times$}：联轴器每半转变向一次 \\
角度不对中 & 1 次 & \textbf{$1\times+2\times$}：椭圆涡动 \\
复合不对中 & — & $1\times\sim4\times$：两者叠加 \\
\bottomrule
\end{tabular}
\end{table}

\textbf{关键认识}：这是一个\textbf{参数激励}系统。
刚度随时间周期变化，响应频谱必然含谐波；
也就是说，不对中的高次谐波\emph{并非}非线性失真所致，
而是\textbf{线性时变系统}的固有结果
——周期系数微分方程的 Floquet 解本身就含各阶谐波。

\subsection{文献}

\begin{table}[h]
\centering
\caption{不对中建模的主要文献（按被引排序）}
\small
\begin{tabular}{@{}p{6.1cm}p{3.4cm}r@{}}
\toprule
文献 & 出处与分量 & 被引 \\
\midrule
Jalan \& Mohanty (2009), \emph{Model based fault diagnosis of a rotor--bearing system for misalignment and unbalance}
& \textbf{JSV} 327:604--622（故障诊断权威期刊，Q1） & \textbf{228} \\
Xu \& Marangoni (1994), Part I, \emph{Theoretical model and analysis}\(\dagger\)
& \textbf{JSV} 176(5):663--679 & \textbf{209} \\
Patel \& Darpe (2009), \emph{Experimental investigations on vibration response of misaligned rotors}\(\dagger\)
& \textbf{MSSP} 23(7):2236--2252（\textbf{信号处理顶刊}） & \textbf{202} \\
Patel \& Darpe (2009), \emph{Vibration response of misaligned rotors}\(\dagger\)
& \textbf{JSV} 325(3):609--628 & \textbf{149} \\
Lees (2007), \emph{Misalignment in rigidly coupled rotors}\(\dagger\)
& \textbf{JSV} 305(1--2):261--271 & \textbf{124} \\
Al-Hussain (2003), \emph{Dynamic stability of two rigid rotors connected by a flexible coupling with angular misalignment}\(\dagger\)
& \textbf{JSV} 266(2):217--234 & 102 \\
Xu \& Marangoni (1994), Part II, \emph{Experimental validation}\(\dagger\)
& \textbf{JSV} 176(5):681--691 & 99 \\
\bottomrule
\end{tabular}
\end{table}

\noindent
其中 Xu \& Marangoni (1994) 是与本课题最贴近的一篇：
它建立的是\textbf{电机--柔性联轴器--转子}整体系统的理论模型，
且\textbf{同时包含不对中与不平衡}两种故障，
这正对应三合一故障中机械部分的建模需求。

\section{傅里叶论证：从周期性到正弦叠加}
\label{sec:fourier}

设故障激励 $g(\theta)$ 是转角 $\theta=\omega_{\rm rot}t$ 的周期函数，
周期为 $2\pi$（或 $2\pi/k$）。则其傅里叶展开为
\begin{equation}
g(\omega_{\rm rot}t)=\sum_{i=1}^{\infty}\left[a_i\cos(i\omega_{\rm rot}t)+b_i\sin(i\omega_{\rm rot}t)\right]
=\sum_{i}p_i\sin(i\omega_{\rm rot}t+\varphi_i).
\end{equation}

由于工程中能量集中在低阶，截断至 $i=n<\infty$ 即可，
即得式 (1) 所定义的集合 $\Phi$。

值得注意的是，$\Phi$ 允许 $p_i(t)$ 为\textbf{多项式}，
恰好覆盖故障幅值随转速、负载慢变的情形：
不平衡幅值 $me\omega^2$ 在变速过程中即为慢变，
落在 $p_i(t)$ 的描述能力之内。

\section{机械故障如何进入转矩通道}

RCESO 的被控对象是转速环，故障只能通过\textbf{负载转矩} $T_L(t)$ 进入控制回路。
因此必须说明：离心力本身不是转矩，它是如何转化为转矩波动的。

\begin{table}[h]
\centering
\caption{两类故障进入转矩通道的物理路径}
\begin{tabular}{@{}p{2.6cm}p{8.4cm}c@{}}
\toprule
故障 & 进入转矩的物理路径 & 阶次 \\
\midrule
不平衡 &
① 径向力调制轴承摩擦，形成阻力矩波动；\newline
② \textbf{气隙偏心}引起不平衡磁拉力（UMP），产生附加转矩谐波 &
$1\times$ \\
不对中 &
联轴器直接传递\textbf{周期性扭矩}（SCS 与 ACS 的角向效应） &
$1\times,2\times$ \\
\bottomrule
\end{tabular}
\end{table}

不平衡这一路径\textbf{必须借助电磁耦合}才能成立，不能直接断言“离心力产生转矩”。
相关文献：

\begin{table}[h]
\centering
\caption{电机侧耦合的文献}
\small
\begin{tabular}{@{}p{8.2cm}p{3.6cm}r@{}}
\toprule
文献 & 出处 & 被引 \\
\midrule
Liu et al. (2022), \emph{Four degrees-of-freedom rotor considering combined eccentricity and electromagnetic effects}
& \textbf{MSSP} 169:108740（\textbf{顶刊}） & 29 \\
Song et al. (2016), \emph{Cogging torque and unbalanced magnetic pull due to simultaneous dynamic and static eccentricities}
& \textbf{IEEE Trans. Magnetics} & — \\
\bottomrule
\end{tabular}
\end{table}

\section{等效载荷与分离辨识（方法层补充）}

除上述机理文献外，故障诊断领域另有一条以\textbf{等效载荷}为核心的方法论线索，
其思想与 EID 高度一致：不重建故障本体，只辨识其等效作用。

\begin{table}[h]
\centering
\caption{等效载荷方法论文献}
\small
\begin{tabular}{@{}p{8.2cm}p{3.4cm}r@{}}
\toprule
文献 & 出处 & 被引 \\
\midrule
Bachschmid, Pennacchi, Vania (2002), \emph{Identification of multiple faults in rotor systems}\(\dagger\)
& \textbf{JSV} 254(2):327--366 & \textbf{162} \\
Kumar \& Tiwari (2021), \emph{Unbalanced and misaligned flexible rotor system}\(\dagger\)
& \textbf{MSSP} 152:107454 & 67 \\
Tiwari \& Kumar (2022), \emph{Virtual trial misalignment approach}\(\dagger\)
& \textbf{MSSP} 167:108540 & 30 \\
\bottomrule
\end{tabular}
\end{table}

\section{开放获取的补充文献}

上节与 \S\ref{sec:mis} 中的多篇为 Elsevier 订阅文献，
经 Semantic Scholar 与 Unpaywall 双渠道核实\textbf{均无合法免费版本}。
为使论证链在无法获取全文时仍然成立，另选取以下 CC-BY 文献作为补充，
已随本文归档于 \texttt{\_refs/} 目录。

\begin{table}[h]
\centering
\caption{开放获取补充文献（非顶刊，仅作可验证材料）}
\small
\begin{tabular}{@{}p{6.4cm}p{4.6cm}@{}}
\toprule
文献 & 覆盖内容 \\
\midrule
Ma, Liu, Li, Zhang, Han (2026), \emph{Nonlinear vibration responses of a rigid--flexible coupled dual-rotor system under coexisting unbalance, misalignment, and rub-impact faults}, \textbf{Machines} 14(11):1337
& \textbf{不平衡 + 不对中共存}的建模与实验，最接近 Xu \& Marangoni 的定位 \\
Xing, Zhang, Wang, Zhang (2024), \emph{The effect of support misalignment on vibration characteristics of aero-engine rotor systems}, \textbf{Machines} 12(10):669
& 不对中\textbf{解析激励模型}的建立 \\
Jin \& Liu (2022), \emph{Rod--disk rotor with inner misalignment}, \textbf{Machines} 10(5):355
& 内部不对中与伴生不平衡 \\
\bottomrule
\end{tabular}
\end{table}

\noindent
\textbf{期刊分量说明}：\emph{Machines} 为 MDPI 开放获取期刊，
学术分量明显低于 IEEE TIE、MSSP 与 JSV，
\textbf{不宜作为主要论据}；此处仅用于在订阅文献不可得时提供可核验的物理描述。

\section{汇总：论证链与文献的对应}

\begin{longtable}{@{}p{6.6cm}p{6.4cm}@{}}
\caption{论证链与文献对应表}\\
\toprule
论证环节 & 支撑文献 \\
\midrule
\endfirsthead
\toprule
论证环节 & 支撑文献 \\
\midrule
\endhead
\bottomrule
\endfoot
等效扰动可折至输入端（存在性） & She et al. (2008), \textbf{IEEE TIE}, 574 次 \\
EID 与 GESO 的关系（方法选型） & She et al. (2023), \textbf{IEEE/CAA JAS}, 79 次 \\
不平衡为严格 $1\times$ & Friswell et al. (2010), \textbf{Cambridge UP}, 422 次 \\
不对中的谐波机制（SCS+ACS 时变刚度） & Lees (2007), \textbf{JSV}, 124 次 \\
电机--联轴器--转子，不对中+不平衡 & Xu \& Marangoni (1994), \textbf{JSV}, 209 次 \\
不对中与不平衡的联合诊断 & Jalan \& Mohanty (2009), \textbf{JSV}, 228 次 \\
不对中响应中的高次谐波（实验） & Patel \& Darpe (2009), \textbf{MSSP}, 202 次 \\
故障即等效载荷（方法论） & Bachschmid et al. (2002), \textbf{JSV}, 162 次 \\
多故障分离辨识 & Kumar \& Tiwari (2021), \textbf{MSSP}, 67 次 \\
偏心 → 电磁力 → 扰动转矩 & Liu et al. (2022), \textbf{MSSP}, 29 次 \\
共存故障的建模与实验（可获取） & Ma et al. (2026), \textbf{Machines}（CC-BY） \\
\end{longtable}

\section{结论}

\begin{enumerate}[label=(\arabic*)]
\item 质量不平衡与不对中之所以可等效为有限个正弦之和，
\textbf{根源在于旋转机械故障的几何周期性}；
正弦形式是傅里叶分解的结果，不是建模假设。

\item 两类故障的阶次特征不同，但均可归入同一扰动集：
不平衡严格为 $1\times$（离心力 $me\omega^2e^{j\omega t}$）；
不对中因联轴器刚度呈 $K_{\rm SCS}+K_{\rm ACS}(\theta)$ 的时变特性而生成 $1\times,2\times,\dots$ 谐波。

\item 不对中的高次谐波是\textbf{线性时变系统}（参数激励）的固有结果，
而非非线性失真，这一点在建模时容易被误解。

\item 故障进入转速环的通道是\textbf{负载转矩}；
不平衡须经由电磁耦合（气隙偏心、UMP）方能产生转矩谐波，
不对中则通过联轴器直接传递周期性扭矩。

\item 本文所列主干文献以 IEEE TIE、MSSP、JSV 与 Cambridge 专著为主，
均为相关领域顶级或权威出处；
开放获取文献仅作可核验的补充材料，不承担主要论据。
\end{enumerate}

\end{document}
